\scanpage{42}
% archive foliation: 41
On trouve donc un foncteur pl. fid
\[
C^\circ\lhook\joinrel\xrightarrow{\deleted{\unsure{\beta}}\alpha^\bullet}\mathrm{Ab}^\bullet_k\deleted{\unsure{*}},
\]
qui associe à tout \deleted{\unsure{M_*}} $L_*$ associe le complexe
de cochaînes $\operatorname{Hom}_C(L_*,C^\bullet_*)$
\diagram{\begin{tikzcd}[column sep=small]
0\arrow[r]&\operatorname{Hom}(L_*,C^0_*)\arrow[r]&\operatorname{Hom}(L_*,C^1_*)\arrow[r]&\cdots
\end{tikzcd}}
et dont la restriction à $A$ est donnée par
\[
\alpha^\bullet\deleted{\unsure{\beta}}(C^i_*)=\left(0\longrightarrow0\longrightarrow\cdots\ 0\longrightarrow
\underset{\text{degré }i}{k}\xrightarrow{\mathrm{id}}
\underset{\text{degré }i+1}{k}\longrightarrow0\longrightarrow\cdots\right)
=\deleted{\unsure{K}}\inserted{C'_i},
\]
et par la loi
\diagram{\begin{tikzcd}[row sep=small,column sep=large]
&\alpha^\bullet(C^i_*)\arrow[d,equal]&\alpha(C^{i+1}_*)\arrow[l,"\alpha^\bullet(d^i)"']\arrow[d,equal]\\
&\unsure{(C'_i)^\bullet|K'_i}\arrow[r]&\unsure{(C'_{i+1})^\bullet|K'_{i+1}}\\
&\vdots&\vdots\\
&0\arrow[d]&0\arrow[d]\\
&k\arrow[d]&0\arrow[l]\arrow[d]\\
i+1&k\arrow[d]&k\arrow[l,leftrightarrow,"\mathrm{id}"']\arrow[d]\\
i+2&0\arrow[d]&k\arrow[l]\arrow[d]\\
&0&0\\
&\vdots&\vdots
\end{tikzcd}}
\[
k\longrightarrow0
\]
\margin{L'image essentielle \inserted{d'$\alpha^\bullet$} de $C^\circ$
est la catégor. des $K^\bullet\in\operatorname{Ob}\mathrm{Ab}_k^\bullet$
qui sont libres de t.f.
en ttes dimensions, et
qui sont à \unsure{ch.} nulle
(on admettra que tt facteur
direct d'un libre de t.f.
est libre).}

\margin{\unsure{Commencer} par Ex 2}
\uline{Ex 2} Faisons $B=\mathrm{Ab}_k^\circ$,
qui est $k$-linéaire (on suppose
$k$ commutatif), et pour $\forall L_*\in\operatorname{Ob}C$, on consi-
dère le foncteur controv. qu'il représente, \unsure{interprété}
comme foncteur $k$-lin.
\[
C^\circ\longrightarrow\deleted{\unsure{\mathrm{Ab}^\circ_k}}\mathrm{Ab}_k^{\unsure{\circ|\bullet}}:
\quad M_*\longmapsto\operatorname{Hom}_C(M_*,L_*)
\]
on trouve ainsi un foncteur pl. fid (pour la
même raison)
\diagram{\begin{tikzcd}[column sep=small,row sep=large]
C\arrow[r,hook]&\underline{\operatorname{Hom}}_{k\text{-}lin}(C^\circ,\mathrm{Ab}_k^{\unsure{\circ|\bullet}})\arrow[r,"\simeq"]\arrow[d,"\simeq"']&
\underline{\operatorname{Hom}}_{k\text{-}lin}(C,\mathrm{Ab}_k^\circ)^\circ\arrow[d,"\simeq"]\\
&\underline{\operatorname{Hom}}_{k\text{-}lin}(A^\circ,\mathrm{Ab}_k^{\unsure{\circ|\bullet}})\arrow[r,"\simeq"]\arrow[d]&
\underline{\operatorname{Hom}}_{k\text{-}lin}(A,\mathrm{Ab}_k^\circ)^\circ\arrow[d,"\simeq"]\\
&\unsure{\{(\mathrm{Ab}_k^\circ)_*\simeq\deleted{\mathrm{Ab}_{k*}}^{*}\}^{\circ}}&
(\mathrm{Ab}_k^{\circ\unsure{*|\bullet}})^\circ
\end{tikzcd}}
on trouve un foncteur pl. fid.
\[
C^{\unsure{\circ}}\lhook\joinrel\xrightarrow{\beta}\mathrm{Ab}_{k*}.
\]
